\section{Method}
\label{sec:method}

This section develops \ours{} following the paper's central principle: derive the target
with policy optimization, then train it with self-distillation. We first show that vanilla OPSD
is precisely the $\beta=1$ member of a broader KL-regularized policy-optimization family,
making the strength of reference regularization explicit and controllable. We then derive the
family's optimal policy as a geometric interpolation between the reference policy and the
privileged teacher. Because direct RL optimization is costly and the exact sequence-level
target is intractable, we return to efficient distillation by realizing the target through
scheduled token-level logit interpolation. Finally, we derive return-to-go credit assignment
to align the practical token-level update with the underlying sequence-level objective.
Appendix~\ref{app:notation} provides a consolidated reference for the notation used
throughout the paper.

\subsection{Vanilla OPSD as \texorpdfstring{$\beta=1$}{beta=1} Policy Optimization}
\label{sec:beta_opsd}

Let $\pi_\theta(y\mid x)$ denote the student policy, $p_T(y\mid x,c)$ the privileged teacher,
and $\pi_{\mathrm{ref}}(y\mid x)$ a reference policy, such as the initial student or a
stop-gradient copy of the current student. We start from the standard KL-regularized RL
objective
\begin{equation}
    \max_{\theta}
    \;
    \mathbb{E}_{y\sim\pi_\theta(\cdot\mid x)}
    \left[
    R(y;x,c)
    \right]
    -
    \beta
    D_{\mathrm{KL}}
    \left(
    \pi_\theta(\cdot\mid x)
    \,\middle\|\,
    \pi_{\mathrm{ref}}(\cdot\mid x)
    \right),
    \label{eq:kl_regularized_rl}
\end{equation}
where $\beta>0$ controls how strongly the learned policy is constrained to remain close to
the reference.

To embed standard on-policy self-distillation within this policy-optimization family, we
choose the reward to be the teacher-to-reference log-ratio:
\begin{equation}
    R(y;x,c)
    =
    \log
    \frac{
    p_T(y\mid x,c)
    }{
    \pi_{\mathrm{ref}}(y\mid x)
    }.
    \label{eq:teacher_log_ratio_reward}
\end{equation}

This reward is high when a trajectory is more likely under the privileged teacher than under
the reference policy. Substituting Eq.~\ref{eq:teacher_log_ratio_reward} into
Eq.~\ref{eq:kl_regularized_rl} gives the family
\begin{equation}
    \mathcal{J}_{\beta}(\theta)
    =
    \mathbb{E}_{y\sim\pi_\theta(\cdot\mid x)}
    \left[
    \log
    \frac{
    p_T(y\mid x,c)
    }{
    \pi_{\mathrm{ref}}(y\mid x)
    }
    \right]
    -
    \beta
    D_{\mathrm{KL}}
    \left(
    \pi_\theta(\cdot\mid x)
    \,\middle\|\,
    \pi_{\mathrm{ref}}(\cdot\mid x)
    \right).
    \label{eq:beta_opsd_objective}
\end{equation}
We call this objective \emph{$\beta$-OPSD}. It generalizes standard on-policy
self-distillation by treating $\beta$ as a variable regularization parameter rather than
implicitly fixing it at one.
The log-ratio reward favors trajectories preferred by the privileged teacher over the
reference policy, while the KL term penalizes movement away from the reference; $\beta$
controls the trade-off between these two forces.

\begin{proposition}[OPSD is the $\beta=1$ case]
\label{prop:opsd_beta_one}
When $\beta=1$, maximizing Eq.~\ref{eq:beta_opsd_objective} is equivalent to minimizing
the standard sequence-level reverse KL,
\begin{equation}
    D_{\mathrm{KL}}
    \left(
    \pi_\theta(\cdot\mid x)
    \,\middle\|\,
    p_T(\cdot\mid x,c)
    \right).
    \label{eq:standard_opsd_kl}
\end{equation}
\end{proposition}

At $\beta=1$, the reference-policy terms in the log-ratio reward and KL penalty cancel,
leaving exactly the negative reverse KL from the student to the privileged teacher. Thus,
standard OPSD is not separate from policy optimization; it is one member of this
KL-regularized family. For $\beta>1$, the remaining reference penalty makes the update
more conservative, while decreasing $\beta$ toward $1$ moves the objective more
aggressively toward the teacher. The proof is given in Appendix~\ref{app:regularized_rl}.

\subsection{Variable \texorpdfstring{$\beta$}{beta} Defines a Path of Optimal Targets}
\label{sec:optimal_policy_interpolant}

Once vanilla OPSD is recognized as the $\beta=1$ case, the next question is what target
each value of $\beta$ defines. For any fixed regularization coefficient, the optimal policy of
$\beta$-OPSD can be written in closed form.

\begin{proposition}[Optimal policy of $\beta$-OPSD]
\label{prop:geometric_optimal_policy}
Assuming compatible support, the maximizer of Eq.~\ref{eq:beta_opsd_objective} is
\begin{equation}
    \pi_{\beta}^{\star}(y\mid x,c)
    =
    \frac{
    \pi_{\mathrm{ref}}(y\mid x)^{1-\frac{1}{\beta}}
    p_T(y\mid x,c)^{\frac{1}{\beta}}
    }{
    Z_{\beta}(x,c)
    },
    \label{eq:sequence_geometric_mixture}
\end{equation}
where $Z_{\beta}(x,c)$ normalizes over full trajectories.
\end{proposition}

The proof is in Appendix~\ref{app:geometric-optimal-policy}. The solution is a geometric
interpolation between the reference policy and the privileged teacher. For $\beta\ge 1$, the
teacher weight $\frac{1}{\beta}$ lies in $[0,1]$: $\beta=1$ recovers the teacher endpoint,
while $\beta\to\infty$ approaches the reference endpoint. Thus, the teacher is no longer
the unique target. Each value of $\beta$ selects a different optimal policy along a principled
reference-to-teacher path.

\paragraph{Scheduling the interpolation.}
The theory above defines a family of targets for fixed $\beta$. In practice, we exploit this
new degree of freedom by varying the coefficient over training. We write the coefficient at
step $k$ as $\beta_k$; its inverse directly determines the teacher weight. We require
\begin{equation}
    \frac{1}{\beta_k}\in[0,1],
    \qquad
    \text{equivalently } \beta_k\ge 1,
\end{equation}
so that the target remains on the path between the reference and the teacher. A larger
$\frac{1}{\beta_k}$ gives a more teacher-like target, while a smaller
$\frac{1}{\beta_k}$ keeps the target closer to the reference policy.

In practice, we use a bounded linear schedule for the teacher weight:
\begin{equation}
    \frac{1}{\beta_k}
    =
    w_{\mathrm{start}}
    +
    \left(
    w_{\mathrm{end}}-w_{\mathrm{start}}
    \right)
    \frac{k}{K-1},
    \qquad
    k=0,\ldots,K-1.
    \label{eq:bounded_beta_schedule}
\end{equation}
Equivalently,
\begin{equation}
    \beta_k
    =
    \left[
    w_{\mathrm{start}}
    +
    \left(
    w_{\mathrm{end}}-w_{\mathrm{start}}
    \right)
    \frac{k}{K-1}
    \right]^{-1}.
\end{equation}
When $w_{\mathrm{start}}<w_{\mathrm{end}}$, early targets remain closer to the
reference policy, while later targets incorporate stronger teacher guidance. This smooth
curriculum avoids an abrupt projection onto the privileged teacher and provides a direct
mechanism for reducing the brittleness of vanilla OPSD.

\paragraph{Local logit realization.}
Although Eq.~\ref{eq:sequence_geometric_mixture} gives the optimal trajectory-level
distribution, it is not directly usable for language models because the normalizer
$Z_{\beta}(x,c)$ ranges over all possible completions. We therefore approximate the
optimal policy locally at each decoding prefix.

Let
\begin{equation}
    h_t=(x,y_{<t})
\end{equation}
be the prefix at token $t$. Let $z_{\mathrm{ref}}(\cdot\mid h_t)$ denote the reference
or student-side logits, and let $z_T(\cdot\mid h_t,c)$ denote the privileged teacher logits.
Given the scheduled coefficient $\beta_k$, we define the local logit-interpolant target as
\begin{equation}
    \tilde p_{\beta_k}(\cdot\mid h_t,c)
    =
    \operatorname{softmax}
    \left(
    \left(1-\frac{1}{\beta_k}\right)
    z_{\mathrm{ref}}(\cdot\mid h_t)
    +
    \frac{1}{\beta_k}
    z_T(\cdot\mid h_t,c)
    \right).
    \label{eq:local_logit_interpolant}
\end{equation}
At each prefix, this distribution is exactly the normalized geometric interpolation of the
local reference and teacher distributions: multiplying probabilities geometrically is
equivalent, up to normalization, to linearly interpolating their logits.

The resulting autoregressive interpolant target is
\begin{equation}
    \tilde p_{\beta_k}(y\mid x,c)
    =
    \prod_{t=1}^{T}
    \tilde p_{\beta_k}(y_t\mid x,y_{<t},c).
    \label{eq:sequence_logit_interpolant}
\end{equation}
Because normalization is applied separately at each prefix, the resulting autoregressive
distribution is a tractable local approximation to the sequence-level optimum in
Eq.~\ref{eq:sequence_geometric_mixture}. This is the target used by the practical
distillation objective below.

\subsection{From Policy Optimization Back to Distillation}
%\subsection{Policy-Optimal Targets via Efficient Self-Distillation}
\label{sec:logit_interpolant_distillation}

The $\beta$-OPSD objective in Eq.~\ref{eq:beta_opsd_objective} can in principle be
optimized with KL-regularized RL algorithms. Doing so, however, would introduce additional
RL machinery, such as reward-weighted policy updates and advantage estimation, together
with higher-variance gradients. The analysis above provides a cheaper route: the policy
objective specifies an optimal sequence-level target, and
Eq.~\ref{eq:sequence_logit_interpolant} gives its tractable local approximation.

This is the main ``back to distillation'' step. Rather than directly maximizing the RL
objective, at training step $k$ we train the student to match the RL-derived interpolant
target $\tilde p_{\beta_k}$ by minimizing
\begin{equation}
    \mathcal{L}_{\beta_k}(\theta)
    =
    D_{\mathrm{KL}}
    \left(
    \pi_\theta(\cdot\mid x)
    \,\middle\|\,
    \tilde p_{\beta_k}(\cdot\mid x,c)
    \right).
    \label{eq:interpolant_distillation_objective}
\end{equation}
Equivalently,
\begin{equation}
    \mathcal{L}_{\beta_k}(\theta)
    =
    \mathbb{E}_{y\sim\pi_\theta(\cdot\mid x)}
    \left[
    \log
    \frac{
    \pi_\theta(y\mid x)
    }{
    \tilde p_{\beta_k}(y\mid x,c)
    }
    \right].
\end{equation}
Using the autoregressive factorization in Eq.~\ref{eq:sequence_logit_interpolant}, this
becomes
\begin{equation}
    \mathcal{L}_{\beta_k}(\theta)
    =
    \mathbb{E}_{y\sim\pi_\theta(\cdot\mid x)}
    \left[
    \sum_{t=1}^{T}
    \rho_t^{\beta_k}(y)
    \right],
    \label{eq:interpolant_kl_decomposition}
\end{equation}
where
\begin{equation}
    \rho_t^{\beta_k}(y)
    =
    \log \pi_\theta(y_t\mid x,y_{<t})
    -
    \log \tilde p_{\beta_k}(y_t\mid x,y_{<t},c).
    \label{eq:interpolant_log_ratio}
\end{equation}

This objective has the same form as vanilla OPSD, except that the fixed teacher
$p_T$ is replaced by the scheduled interpolant target $\tilde p_{\beta_k}$. When
$\beta_k=1$, the interpolant target reduces to the privileged teacher, and
Eq.~\ref{eq:interpolant_distillation_objective} recovers standard OPSD. For larger
$\beta_k$, the target remains closer to the reference policy, yielding a more conservative
distillation objective.

The interpolant target is efficient to compute. Once the reference/student logits and
teacher logits are available, constructing $\tilde p_{\beta_k}$ only requires an elementwise
weighted sum of the two logit vectors followed by a softmax, as in
Eq.~\ref{eq:local_logit_interpolant}. Importantly, this avoids the intractable
sequence-level normalizer $Z_\beta(x,c)$ in Eq.~\ref{eq:sequence_geometric_mixture}.
Thus, policy optimization specifies where the student should move, while token-level
distillation provides an inexpensive way to get there.

\subsection{Practical \texorpdfstring{$\beta$}{beta}-OPSD with Return-to-Go Credit}
\label{sec:practical_slide_objective}
We now describe the practical training loss. At each update, we sample trajectories
on-policy from the current student:
\begin{equation}
    y\sim\pi_\theta(\cdot\mid x).
\end{equation}
Along the sampled trajectory, we evaluate the token-level mismatch between the student and
the interpolant target:
\begin{equation}
    \rho_t^{\beta_k}(y)
    =
    \log \pi_\theta(y_t\mid x,y_{<t})
    -
    \log \tilde p_{\beta_k}(y_t\mid x,y_{<t},c).
    \label{eq:practical_interpolant_log_ratio}
\end{equation}
The scalar estimator $\sum_t \rho_t^{\beta_k}(y)$ is an unbiased Monte Carlo estimate of
the sequence-level reverse KL in Eq.~\ref{eq:interpolant_distillation_objective}. However,
using each token's instantaneous log-ratio as its training weight gives a myopic gradient:
it ignores how an earlier token affects future prefixes and future target mismatch.

Therefore, $\beta$-OPSD uses a return-to-go weight:
\begin{equation}
    G_{t,\gamma}^{\beta_k}(y)
    =
    \sum_{s=t}^{T}
    \gamma^{s-t}
    \rho_s^{\beta_k}(y),
    \qquad
    \gamma\in[0,1].
    \label{eq:slide_return_to_go}
\end{equation}
When $\gamma=1$, this recovers the exact sequence-level gradient estimator; in practice,
we use $\gamma<1$ to control the magnitude of the return on long generations. The proof is in Appendix~\ref{app:onpolicy-lookahead-gradient}.

The practical $\beta$-OPSD loss is
\begin{equation}
    \mathcal{L}_{\beta}(\theta;y)
    =
    \frac{1}{T}
    \sum_{t=1}^{T}
    \operatorname{sg}
    \left(
    G_{t,\gamma}^{\beta_k}(y)
    \right)
    \log \pi_\theta(y_t\mid x,y_{<t}),
    \label{eq:slide_loss}
\end{equation}
where $\operatorname{sg}(\cdot)$ denotes stop-gradient. 
The interpolant logits and
return-to-go weights are detached; the only gradient path is through the student
log-probability. Thus, $\beta$-OPSD differs from vanilla OPSD in two simple ways:
\begin{enumerate}
    \item it replaces the fixed teacher target with a scheduled logit-interpolant target, and
    \item it replaces local token weights with return-to-go weights for sequence-level credit
    assignment.
\end{enumerate}

The first modification is the main conceptual contribution: it follows from the
$\beta$-OPSD optimal policy. The second is a practical sequence-level estimator that makes
the distillation loss faithful to the trajectory objective.
